\scanpage{25}
\uline{Orientations}

a) Vectoriel réel $E$ de dim finie $n$. La donnée de $E$ équivaut à celle du torseur à dr.
\[
 P=\operatorname{Isom}(\mathbb R^n,E)\simeq\operatorname{bases}(E)
\]
sous
\[
 G=\operatorname{Gl}(n,\mathbb R).
\]

On a un hom
\diagram{\begin{tikzcd}[ampersand replacement=\&,column sep=small,row sep=small]
 G=\operatorname{Gl}(n,\deleted{\unsure{R|k}}\mathbb R)\arrow[rr,"s"]\arrow[dr,"\det"']
 \& {} \& \{\pm1\} \& u\longmapsto\deleted{\det}\,\operatorname{sgn}\det u\\
 {} \& \mathbb R^{\ast}\arrow[ur,"\operatorname{sgn}"'] \& {}
 \& \text{noyau }\operatorname{Gl}^+(n,\mathbb R)\,,
\end{tikzcd}}
et
\[
 \omega_E=P\times_G\{\pm1,s\}\simeq P/\operatorname{Gl}^+(n,R)
\]
est l'ens des orientations de $E$.

b) Module \inserted{\unsure{quadratique|quadratiques}} libre \deleted{\unsure{quadr.|point}} \inserted{$E$} sur un anneau comm. $k$ \unsure{(où|où} $2$ est inversible [$\exists$ base orthonormale, $\deleted{\unsure{Q(e_i,e|Q(e_i,e_j)}} e_i^2=1$, $e_i e_j=0$ si $i\ne j$]

La donnée de $E$ de rang $n$ équivaut à celle du tors à droite
\[
 P=\operatorname{Isom}_{k\text{-quadr.}}(k^n,E)\simeq\text{bases orthon.}(E)
\]
sous
\[
 G=\operatorname{O}(n,k)=\operatorname{Aut}_{k\text{-quadr}}(k^n,Q_n)\subset\operatorname{Gl}(n\unsure{.|,}k)
\]

\unsure{Or|On} ici
\[
 G=\operatorname{O}(n,k)\xrightarrow{\substack{\det=s\\\unsure{\deleted{\operatorname{sgn}}}}}\{\pm1\}\qquad(\text{noyau }\operatorname{SO}(n,k))
\]
et l'ens des orientations de $E$ est (par définition)
\[
 \omega_E=P\times_G(\pm1,s)\simeq P/\operatorname{SO}(n,k).
\]

\uline{Compatibilité} \uline{entre a) et b)} Résulte de la transitivité de l'opération d'ext. du groupe d'opérateurs dans un torseur, et du diagramme comm.
\diagram{\begin{tikzcd}[ampersand replacement=\&,column sep=large]
 \operatorname{O}(n,k)\arrow[r,hook,"\mathrm{incl}"]\arrow[rr,bend right=19,"\det"] \& \operatorname{Gl}(n,k)\arrow[r,"\operatorname{sgn}\det"] \& \{\pm1\}
\end{tikzcd}}
% archive sheet number: 24
